% !TEX program = pdflatex
% Hugging Face Deep RL Course - Unit 4 (complete MDX study edition).
% Compiled with latexmk -pdf from the notes/ directory; uses existing preamble.tex.
\PassOptionsToPackage{T2A}{fontenc}
\PassOptionsToPackage{russian}{babel}
\documentclass[11pt,a4paper]{article}

% Shared typography/layout for Hugging Face Deep RL Course notes.
% Compile with pdfLaTeX / latexmk -pdf with notes/ as the working directory.
% e.g. cd notes && latexmk -pdf unit1.tex

\usepackage[T1]{fontenc}
\usepackage[utf8]{inputenc}
\usepackage[english]{babel}
\usepackage{lmodern}
\usepackage{textcomp}
\usepackage{geometry}
\geometry{a4paper,margin=22mm,headheight=14pt,headsep=8mm,footskip=12mm}

\usepackage{amsmath,amssymb,mathtools}
\usepackage{graphicx,caption,needspace,placeins}
\usepackage{array,booktabs,tabularx}
\usepackage{xcolor}
\usepackage{microtype}
\usepackage{enumitem}
\usepackage{listings}
\usepackage[most]{tcolorbox}
\usepackage{titlesec}
\usepackage{fancyhdr}
\usepackage{hyperref}
\usepackage{bookmark}

\definecolor{ink}{HTML}{243041}
\definecolor{accent}{HTML}{38598C}
\definecolor{secondary}{HTML}{60768A}
\definecolor{codebg}{HTML}{F4F6F8}
\definecolor{codeborder}{HTML}{E3E8EE}
\definecolor{notebg}{HTML}{F3F7FC}

\hypersetup{
  colorlinks=true,
  linkcolor=accent,
  urlcolor=accent,
  citecolor=accent,
  pdfborder={0 0 0}
}
\urlstyle{same}

\setlength{\parindent}{0pt}
\setlength{\parskip}{0.38em plus 0.12em minus 0.05em}

\setlist[itemize]{
  topsep=0.35em,
  itemsep=0.13em,
  leftmargin=1.7em
}
\setlist[enumerate]{
  topsep=0.35em,
  itemsep=0.13em,
  leftmargin=1.8em
}

\setcounter{tocdepth}{2}
\setcounter{secnumdepth}{2}

\titleformat{\section}
  {\Large\bfseries\color{ink}}
  {\thesection}{0.7em}{}

\titleformat{\subsection}
  {\large\bfseries\color{accent}}
  {\thesubsection}{0.7em}{}

\titleformat{\subsubsection}
  {\normalsize\bfseries\color{ink}}
  {\thesubsubsection}{0.6em}{}

\titlespacing*{\section}
  {0pt}{2.2ex plus 0.7ex}{0.9ex}
\titlespacing*{\subsection}
  {0pt}{1.7ex plus 0.4ex}{0.6ex}
\titlespacing*{\subsubsection}
  {0pt}{1.3ex plus 0.3ex}{0.45ex}

\pagestyle{fancy}
\fancyhf{}

\fancyhead[L]{
  \footnotesize\sffamily\color{secondary}
  HF Deep RL Course \enspace/\enspace Unit 1
}
\fancyhead[R]{
  \footnotesize\sffamily\color{secondary}
  \nouppercase{\rightmark}
}
\fancyfoot[C]{
  \footnotesize\color{secondary}\thepage
}
\renewcommand{\headrulewidth}{0.25pt}

\captionsetup{
  font=small,
  labelfont={bf,color=secondary},
  textfont={color=ink},
  justification=centering,
  skip=4pt,
  hypcap=false
}

\lstdefinestyle{hfcode}{
  basicstyle=\ttfamily\footnotesize,
  backgroundcolor=\color{codebg},
  frame=single,
  rulecolor=\color{codeborder},
  numbers=none,
  showstringspaces=false,
  breaklines=true,
  breakatwhitespace=false,
  columns=fullflexible,
  keepspaces=true,
  xleftmargin=0pt,
  xrightmargin=0pt,
  aboveskip=0.6em,
  belowskip=0.6em,
  framexleftmargin=7pt,
  framexrightmargin=7pt,
  framextopmargin=5pt,
  framexbottommargin=5pt,
  upquote=true,
  tabsize=4,
}
\lstset{style=hfcode}

\newcommand{\CourseFigure}[2]{%
  \par\addvspace{0.55em}\noindent%
  \begin{minipage}{\linewidth}%
    \centering
    \includegraphics[
      width=0.86\linewidth,
      height=0.295\textheight,
      keepaspectratio
    ]{#1}\par
    \captionof{figure}{#2}%
  \end{minipage}%
  \par\addvspace{0.55em}%
}

\newtcolorbox{coursenote}{
  enhanced,
  breakable,
  colback=notebg,
  colframe=accent!30,
  boxrule=0.35pt,
  arc=1.2mm,
  left=9pt,
  right=9pt,
  top=6pt,
  bottom=6pt
}

\renewenvironment{quote}
  {\begin{coursenote}}
  {\end{coursenote}}

\emergencystretch=2.5em

\providecommand{\tightlist}{
  \setlength{\itemsep}{0pt}
  \setlength{\parskip}{0pt}
}

\providecommand{\passthrough}[1]{#1}

\Urlmuskip=0mu plus 6mu\relax

\renewcommand{\rmdefault}{cmr}
\renewcommand{\sfdefault}{cmss}
\renewcommand{\ttdefault}{cmtt}
\graphicspath{{../assets/unit4/}}
% Reuse existing shared preamble, override unit-specific running headers.
\fancyhead[L]{\footnotesize\sffamily\color{secondary}Курс Deep RL от Hugging Face \enspace/\enspace Модуль 4}
\fancyhead[R]{\footnotesize\sffamily\color{secondary}Градиент политики}
\title{\textbf{Глубокое обучение с подкреплением}\\[0.35em]\large Модуль 4: градиент политики и REINFORCE}
\author{Курс Deep RL от Hugging Face\\\small Полное учебное издание на основе исходных MDX-файлов}
\date{}
\begin{document}
\selectlanguage{russian}
\fontencoding{T2A}\selectfont
\maketitle
\vspace{-0.8em}
\begin{center}\small\color{secondary}Исходный материал: курс Deep RL от Hugging Face.\end{center}
\begingroup\small\setlength{\parskip}{0pt}\tableofcontents\endgroup
\vspace{0.5em}
\clearpage

\section{Введение}\label{introduction}

\CourseFigure{thumbnail.png}{Обложка модуля 4: градиент политики и REINFORCE \protect\href{https://huggingface.co/datasets/huggingface-deep-rl-course/course-images/resolve/main/en/unit6/thumbnail.png}{[оригинал]}}

В предыдущем модуле мы изучили глубокое Q-обучение (Deep Q-Learning). В этом алгоритме глубокого обучения с подкреплением, основанном на оценке ценности, мы \textbf{использовали глубокую нейронную сеть для аппроксимации различных Q-значений каждого возможного действия в данном состоянии.}

С самого начала курса мы изучали только методы, основанные на оценке ценности, \textbf{где функция ценности служит промежуточным этапом на пути к оптимальной политике.}

\CourseFigure{link-value-policy.jpg}{Связь между оценками ценности действий и выбором политики \protect\href{https://huggingface.co/datasets/huggingface-deep-rl-course/course-images/resolve/main/en/unit3/link-value-policy.jpg}{[оригинал]}}

В методах, основанных на ценности, политика \textbf{\(\pi\) возникает на основе оценок ценности действий: по сути, это функция} (например, жадная политика), которая в заданном состоянии выбирает действие с максимальной ценностью.

В методах, основанных на политике, мы стремимся оптимизировать непосредственно политику \textbf{без промежуточного этапа обучения функции ценности.}

Сегодня мы \textbf{познакомимся с методами, основанными на политике, и изучим их подмножество --- методы градиента политики (policy gradient)}. Затем с нуля реализуем на PyTorch первый алгоритм градиента политики --- Monte Carlo \textbf{REINFORCE}. После этого проверим устойчивость его работы в средах CartPole-v1 и PixelCopter.

Далее ты сможешь самостоятельно дорабатывать эту реализацию и применять её в более сложных средах.

\CourseFigure{envs.png}{Среды CartPole и PixelCopter (статичный кадр анимации) \protect\href{https://huggingface.co/datasets/huggingface-deep-rl-course/course-images/resolve/main/en/unit6/envs.gif}{[оригинал]}}

Приступим!

\section{Что такое методы, основанные на политике?}\label{what-are-the-policy-based-methods}

Основная задача обучения с подкреплением --- \textbf{найти оптимальную политику \(\pi^{*}\), максимизирующую ожидаемую суммарную награду}. В основе обучения с подкреплением лежит \emph{гипотеза вознаграждения} (reward hypothesis): \textbf{любую цель можно выразить как максимизацию ожидаемой суммарной награды.}

Например, в футболе (агентов для которого ты будешь обучать через два модуля) цель --- выиграть матч. На языке обучения с подкреплением её можно выразить как \textbf{максимизацию числа забитых сопернику голов} (когда мяч пересекает линию ворот) и \textbf{минимизацию числа голов, пропущенных в собственные ворота}.

\CourseFigure{soccer.jpg}{Пример футбола для объяснения цели максимизации награды \protect\href{https://huggingface.co/datasets/huggingface-deep-rl-course/course-images/resolve/main/en/unit6/soccer.jpg}{[оригинал]}}

\subsection{Методы на основе ценности, на основе политики и Actor--Critic}\label{value-based-policy-based-and-actor-critic-methods}

В первом модуле мы рассмотрели два подхода к поиску (или, чаще, аппроксимации) этой оптимальной политики \(\pi^{*}\).

\begin{itemize}
\tightlist
\item
  В \emph{методах на основе ценности} мы обучаем функцию ценности.

  \begin{itemize}
  \tightlist
  \item
    Идея состоит в том, что оптимальная функция ценности позволяет получить оптимальную политику \(\pi^{*}\).
  \item
    Наша цель --- \textbf{минимизировать ошибку между предсказанной и целевой ценностью}, чтобы аппроксимировать истинную функцию ценности действий.
  \item
    Политика при этом существует, но задана неявно, поскольку \textbf{непосредственно строится по функции ценности}. Например, в Q-обучении мы использовали \(\varepsilon\)-жадную политику.
  \end{itemize}
\item
  Напротив, в \emph{методах на основе политики} мы напрямую учимся аппроксимировать \(\pi^{*}\), не обучая функцию ценности.

  \begin{itemize}
  \tightlist
  \item
    Основная идея --- \textbf{параметризовать политику}. Например, нейронная сеть \(\pi_\theta\) выдаёт распределение вероятностей действий (стохастическая политика).
  \item
  \end{itemize}
\end{itemize}

\CourseFigure{stochastic_policy.png}{Стохастическая политика выдаёт распределение вероятностей действий \protect\href{https://huggingface.co/datasets/huggingface-deep-rl-course/course-images/resolve/main/en/unit6/stochastic_policy.png}{[оригинал]}}

\begin{itemize}
\tightlist
\item
  Следовательно, наша задача --- \textbf{максимизировать качество параметризованной политики методом градиентного подъёма}.
\item
  Для этого мы изменяем параметры \(\theta\), влияющие на распределение действий в каждом состоянии.
\end{itemize}

\CourseFigure{policy_based.png}{Нейронная сеть политики: преобразование состояния в вероятности действий \protect\href{https://huggingface.co/datasets/huggingface-deep-rl-course/course-images/resolve/main/en/unit6/policy_based.png}{[оригинал]}}

\begin{itemize}
\tightlist
\item
  В следующем модуле мы изучим метод \emph{Actor--Critic}, объединяющий подходы на основе ценности и на основе политики.
\end{itemize}

Таким образом, методы на основе политики позволяют непосредственно оптимизировать \(\pi_\theta\), чтобы получать распределение \(\pi_\theta(a|s)\), обеспечивающее наибольшую суммарную отдачу. Для этого мы определяем целевую функцию \(J(\theta)\), равную ожидаемой суммарной награде, и \textbf{ищем такое значение \(\theta\), при котором она максимальна}.

\subsection{Различие между методами на основе политики и методами градиента политики}\label{the-difference-between-policy-based-and-policy-gradient-methods}

Методы градиента политики, которым посвящён этот модуль, представляют собой подкласс методов на основе политики. Для таких методов оптимизация чаще всего происходит \emph{on-policy}, поскольку на каждом шаге обновления мы используем только данные (траектории), собранные \textbf{самой последней версией политики} \(\pi_\theta\).

Разница между этими подходами \textbf{заключается в способе оптимизации параметров} \(\theta\):

\begin{itemize}
\tightlist
\item
  В \emph{методах на основе политики} мы непосредственно ищем оптимальную политику. Параметры \(\theta\) можно оптимизировать \textbf{косвенно}, максимизируя локальную аппроксимацию целевой функции с помощью восхождения на холм (hill climbing), имитации отжига или эволюционных стратегий.
\item
  В \emph{методах градиента политики}, как подклассе методов на основе политики, мы также непосредственно ищем оптимальную политику. Но параметры \(\theta\) оптимизируются \textbf{напрямую} посредством градиентного подъёма целевой функции \(J(\theta)\).
\end{itemize}

Прежде чем подробно разбирать работу методов градиента политики (целевую функцию, теорему о градиенте политики, градиентный подъём и т. д.), рассмотрим их преимущества и недостатки.

\section{Преимущества и недостатки методов градиента политики}\label{the-advantages-and-disadvantages-of-policy-gradient-methods}

На этом этапе можно спросить: «Но ведь глубокое Q-обучение отлично работает! Зачем использовать градиент политики?» Чтобы ответить, рассмотрим \textbf{преимущества и недостатки методов градиента политики}.

\subsection{Преимущества}\label{advantages}

По сравнению с методами на основе ценности у них есть несколько достоинств. Рассмотрим некоторые из них.

\subsubsection{Простота интеграции}\label{the-simplicity-of-integration}

Мы можем непосредственно оценивать политику, не сохраняя дополнительных данных (ценностей действий).

\subsubsection{Методы градиента политики способны обучать стохастическую политику}\label{policy-gradient-methods-can-learn-a-stochastic-policy}

Методы градиента политики могут \textbf{обучать стохастическую политику, тогда как функции ценности сами по себе такой возможности не дают}.

Отсюда следуют два последствия:

\begin{enumerate}
\def\labelenumi{\arabic{enumi}.}
\item
  Нам \textbf{не требуется вручную задавать компромисс между исследованием среды и использованием найденного (exploration/exploitation)}. Поскольку на выходе получаем распределение вероятностей действий, агент исследует \textbf{пространство состояний, не повторяя постоянно одну и ту же траекторию.}
\item
  Мы также избавляемся от проблемы \textbf{неразличимости состояний по наблюдениям} (perceptual aliasing). Она возникает, когда два состояния выглядят одинаковыми (или действительно совпадают с точки зрения наблюдений), но требуют разных действий.
\end{enumerate}

Рассмотрим пример: у нас есть умный робот-пылесос, который должен собирать пыль и при этом не причинять вреда хомякам.

\CourseFigure{hamster1.jpg}{Стохастические политики: первая иллюстрация с роботом-пылесосом \protect\href{https://huggingface.co/datasets/huggingface-deep-rl-course/course-images/resolve/main/en/unit6/hamster1.jpg}{[оригинал]}}

Наш робот-пылесос умеет определять только расположение стен.

Проблема в том, что \textbf{два отмеченных красным состояния неразличимы для агента: в обоих он воспринимает стену сверху и снизу}.

\CourseFigure{hamster2.jpg}{Стохастические политики: вторая иллюстрация \protect\href{https://huggingface.co/datasets/huggingface-deep-rl-course/course-images/resolve/main/en/unit6/hamster2.jpg}{[оригинал]}}

При детерминированной политике агент в красном состоянии либо всегда идёт вправо, либо всегда влево. \textbf{В обоих случаях он застрянет и так и не доберётся до пыли}.

Если использовать алгоритм обучения с подкреплением на основе ценности, то мы получим \textbf{почти детерминированную политику} («\(\varepsilon\)-жадную стратегию»). В результате агент может \textbf{очень долго искать пыль}.

Напротив, оптимальная стохастическая политика \textbf{случайным образом выбирает движение влево или вправо в красных состояниях}. Благодаря этому \textbf{агент не застрянет и с высокой вероятностью достигнет цели}.

\CourseFigure{hamster3.jpg}{Стохастические политики: третья иллюстрация \protect\href{https://huggingface.co/datasets/huggingface-deep-rl-course/course-images/resolve/main/en/unit6/hamster3.jpg}{[оригинал]}}

\subsubsection{Методы градиента политики эффективнее в пространствах действий высокой размерности и с непрерывными действиями}\label{policy-gradient-methods-are-more-effective-in-high-dimensional-action-spaces-and-continuous-actions-spaces}

Проблема глубокого Q-обучения в том, что для каждого состояния на каждом шаге алгоритм \textbf{предсказывает оценку (максимальную ожидаемую будущую награду) для каждого возможного действия}.

Но что, если возможных действий бесконечно много?

Например, у беспилотного автомобиля в каждом состоянии может быть (практически) бесконечное множество действий: повернуть руль на 15°, 17,2°, 19,4°, посигналить и т. д. \textbf{Нам пришлось бы получать Q-значение для каждого возможного действия}! К тому же \textbf{поиск максимума среди непрерывного множества действий сам по себе представляет задачу оптимизации}!

Вместо этого методы градиента политики выдают \textbf{распределение вероятностей по действиям}.

\subsubsection{Методы градиента политики обладают более благоприятными свойствами сходимости}\label{policy-gradient-methods-have-better-convergence-properties}

В методах на основе ценности мы используем достаточно жёсткий оператор для \textbf{изменения функции ценности: выбираем максимум среди Q-оценок}. Поэтому вероятности выбора действий могут резко меняться даже при сколь угодно малом изменении оценок, если в результате максимальное значение получает другое действие.

Например, если во время обучения лучшим действием было движение влево (Q-значение 0,22), а уже после следующего шага --- вправо (Q-значение для него стало 0,23), то политика резко изменится: теперь агент почти всегда будет двигаться вправо, а не влево.

Напротив, в методах градиента политики предпочтения стохастической политики (вероятности выбора действий) \textbf{изменяются плавно с течением времени}.

\subsection{Недостатки}\label{disadvantages}

Разумеется, методы градиента политики имеют и свои недостатки:

\begin{itemize}
\tightlist
\item
  \textbf{Методы градиента политики нередко сходятся к локальному максимуму вместо глобального оптимума.}
\item
  Градиентный подъём происходит медленно, \textbf{шаг за шагом: обучение может занимать больше времени (неэффективность по выборке)}.
\item
  Оценки градиента политики могут иметь высокую дисперсию. В модуле об Actor--Critic мы разберём причины и способы решения этой проблемы.
\end{itemize}

Если хочешь подробнее узнать о преимуществах и недостатках методов градиента политики, \href{https://youtu.be/y3oqOjHilio}{посмотри это видео}.

\section{Более глубокое знакомство с методами градиента политики}\label{diving-deeper-into-policy-gradient-methods}

\subsection{Общая картина}\label{getting-the-big-picture}

Мы только что узнали, что методы градиента политики ищут параметры \(\theta\), \textbf{максимизирующие ожидаемую отдачу}.

Основная идея: у нас есть \emph{параметризованная стохастическая политика}. В нашем случае нейронная сеть выдаёт распределение вероятностей действий. Вероятность выбора каждого действия также называют \emph{предпочтением действия} (action preference).

Рассмотрим среду CartPole-v1:

\begin{itemize}
\tightlist
\item
  На вход подаётся состояние.
\item
  На выходе получаем распределение вероятностей возможных действий в этом состоянии.
\end{itemize}

\CourseFigure{policy_based.png}{Нейронная сеть политики: преобразование состояния в вероятности действий \protect\href{https://huggingface.co/datasets/huggingface-deep-rl-course/course-images/resolve/main/en/unit6/policy_based.png}{[оригинал]}}

В методах градиента политики мы хотим \textbf{управлять распределением вероятностей действий}, настраивая политику так, чтобы \textbf{хорошие действия (повышающие отдачу) в дальнейшем выбирались чаще}. После каждого взаимодействия агента со средой мы корректируем параметры, увеличивая вероятность выбора удачных действий в будущем.

Но \textbf{как оптимизировать веса с помощью ожидаемой отдачи}?

Идея заключается в том, чтобы \textbf{позволить агенту взаимодействовать со средой в течение эпизода}. Если эпизод завершился успехом, считаем, что выполненные действия были хорошими и впоследствии их стоит выбирать чаще, поскольку они привели к победе.

Следовательно, для каждой пары «состояние--действие» мы хотим увеличить \(P(a|s)\) --- вероятность выбрать данное действие в данном состоянии. А если эпизод завершился неудачей, эту вероятность нужно уменьшить.

Упрощённая схема алгоритма градиента политики выглядит так:

\CourseFigure{pg_bigpicture.jpg}{Общая схема обновления параметров в методе градиента политики \protect\href{https://huggingface.co/datasets/huggingface-deep-rl-course/course-images/resolve/main/en/unit6/pg_bigpicture.jpg}{[оригинал]}}

Теперь, когда общая картина понятна, углубимся в устройство методов градиента политики.

\subsection{Подробный разбор методов градиента политики}\label{diving-deeper-into-policy-gradient-methods-1}

У нас есть стохастическая политика \(\pi\) с параметрами \(\theta\). Получая состояние, эта политика \textbf{выдаёт распределение вероятностей действий}.

\CourseFigure{stochastic_policy.png}{Стохастическая политика выдаёт распределение вероятностей действий \protect\href{https://huggingface.co/datasets/huggingface-deep-rl-course/course-images/resolve/main/en/unit6/stochastic_policy.png}{[оригинал]}}

Здесь \(\pi_\theta(a_t|s_t)\) --- вероятность того, что агент, следуя заданной политике, выберет действие \(a_t\) в состоянии \(s_t\).

\textbf{Но как определить, хороша ли наша политика?} Нужна количественная мера её качества. Для этого введём функцию оценки, или целевую функцию, \(J(\theta)\).

\subsubsection{Целевая функция}\label{the-objective-function}

\emph{Целевая функция} характеризует \textbf{качество работы агента} с учётом траектории (последовательности состояний и действий без учёта наград, в отличие от полного описания эпизода) и возвращает \emph{ожидаемую суммарную награду}.

\CourseFigure{objective.jpg}{Целевая функция градиента политики \protect\href{https://huggingface.co/datasets/huggingface-deep-rl-course/course-images/resolve/main/en/unit6/objective.jpg}{[оригинал]}}

Разберём эту формулу подробнее:

\begin{itemize}
\tightlist
\item
  \emph{Ожидаемая отдача} (или ожидаемая суммарная награда) представляет собой взвешенное среднее всех возможных значений отдачи \(R(\tau)\), причём веса определяются вероятностями \(P(\tau;\theta)\).
\end{itemize}

\CourseFigure{expected_reward.png}{Ожидаемая отдача как целевая функция оптимизации \protect\href{https://huggingface.co/datasets/huggingface-deep-rl-course/course-images/resolve/main/en/unit6/expected_reward.png}{[оригинал]}}

\begin{itemize}
\item
  \(R(\tau)\) --- отдача от произвольной траектории. Чтобы использовать её для вычисления ожидаемой отдачи, нужно умножить её на вероятность соответствующей траектории.
\item
  \(P(\tau;\theta)\) --- вероятность конкретной траектории \(\tau\). Она зависит от \(\theta\), поскольку эти параметры определяют политику выбора действий, а значит, влияют и на посещаемые состояния.
\end{itemize}

\CourseFigure{probability.png}{Вероятность траектории при параметризованной политике \protect\href{https://huggingface.co/datasets/huggingface-deep-rl-course/course-images/resolve/main/en/unit6/probability.png}{[оригинал]}}

\begin{itemize}
\tightlist
\item
  \(J(\theta)\) --- ожидаемая отдача: для её вычисления суммируем по всем траекториям произведения вероятности каждой траектории при заданных \(\theta\) на отдачу этой траектории.
\end{itemize}

Таким образом, наша цель --- максимизировать ожидаемую суммарную награду, найдя параметры \(\theta\), формирующие наилучшие распределения вероятностей действий:

\CourseFigure{max_objective.png}{Максимизация целевой функции по параметрам политики \protect\href{https://huggingface.co/datasets/huggingface-deep-rl-course/course-images/resolve/main/en/unit6/max_objective.png}{[оригинал]}}

\subsection{Градиентный подъём и теорема о градиенте политики}\label{gradient-ascent-and-the-policy-gradient-theorem}

Градиент политики --- задача оптимизации: мы хотим найти значения \(\theta\), при которых целевая функция \(J(\theta)\) максимальна. Поэтому используем \textbf{градиентный подъём} (gradient ascent). Это противоположность \emph{градиентному спуску}: градиент задаёт направление наиболее быстрого увеличения \(J(\theta)\).

Если нужно вспомнить различие между градиентным спуском и подъёмом, прочитай \href{https://www.baeldung.com/cs/gradient-descent-vs-ascent}{эту статью} и \href{https://stats.stackexchange.com/questions/258721/gradient-ascent-vs-gradient-descent-in-logistic-regression}{это объяснение}.

Шаг обновления при градиентном подъёме записывается так:

\(\theta \leftarrow \theta + \alpha *  \nabla_\theta J(\theta)\)

Мы можем многократно применять это обновление, рассчитывая на сходимость \(\theta\) к значению, максимизирующему \(J(\theta)\).

Однако при вычислении производной \(J(\theta)\) возникают две проблемы:

\begin{enumerate}
\def\labelenumi{\arabic{enumi}.}
\item
  Мы не можем вычислить точный градиент целевой функции: для этого нужно рассчитать вероятности всех возможных траекторий, что вычислительно чрезвычайно дорого. Поэтому мы хотим \textbf{оценивать градиент по выборке (собирая некоторое число траекторий)}.
\item
  Есть и другая проблема, которую я объясняю в следующем необязательном разделе. Для дифференцирования целевой функции пришлось бы дифференцировать распределение состояний, связанное с динамикой марковского процесса принятия решений (MDP). Эта динамика определяется средой: она задаёт вероятность перехода в следующее состояние при известном текущем состоянии и действии агента. Но мы можем не знать эту динамику и потому не можем её дифференцировать.
\end{enumerate}

\CourseFigure{probability.png}{Вероятность траектории при параметризованной политике \protect\href{https://huggingface.co/datasets/huggingface-deep-rl-course/course-images/resolve/main/en/unit6/probability.png}{[оригинал]}}

К счастью, существует решение --- \emph{теорема о градиенте политики} (Policy Gradient Theorem). Она позволяет преобразовать выражение для градиента целевой функции так, чтобы не требовалось дифференцировать распределение состояний.

\CourseFigure{policy_gradient_theorem.png}{Теорема о градиенте политики \protect\href{https://huggingface.co/datasets/huggingface-deep-rl-course/course-images/resolve/main/en/unit6/policy_gradient_theorem.png}{[оригинал]}}

Если хочешь понять, как выводится эта формула оценки градиента, перейди к следующему (необязательному) разделу.

\subsection{Алгоритм REINFORCE (Monte Carlo REINFORCE)}\label{the-reinforce-algorithm-monte-carlo-reinforce}

Алгоритм REINFORCE, также называемый методом градиента политики Монте-Карло, --- это алгоритм градиента политики, который \textbf{использует оценку отдачи за целый эпизод для обновления параметров политики} \(\theta\):

В цикле:

\begin{itemize}
\tightlist
\item
  Используем политику \(\pi_\theta\), чтобы собрать эпизод \(\tau\).
\item
  По этому эпизоду оцениваем градиент \(\hat{g} = \nabla_\theta J(\theta)\).
\end{itemize}

\CourseFigure{policy_gradient_one.png}{Обновление REINFORCE по одной выбранной траектории \protect\href{https://huggingface.co/datasets/huggingface-deep-rl-course/course-images/resolve/main/en/unit6/policy_gradient_one.png}{[оригинал]}}

\begin{itemize}
\tightlist
\item
  Обновляем веса политики: \(\theta \leftarrow \theta + \alpha \hat{g}\).
\end{itemize}

Это обновление можно интерпретировать следующим образом:

\begin{itemize}
\item
  \(\nabla_\theta log \pi_\theta(a_t|s_t)\) --- направление \textbf{наиболее быстрого увеличения логарифма вероятности} выбора действия \(a_t\) в состоянии \(s_t\). Этот градиент показывает, \textbf{как нужно изменить веса политики}, чтобы увеличить или уменьшить логарифм вероятности выбора \(a_t\) в \(s_t\).
\item
  \(R(\tau)\) --- функция оценки качества:

  \begin{itemize}
  \tightlist
  \item
    Если отдача велика, она будет \textbf{повышать вероятности} соответствующих пар (состояние, действие).
  \item
    Если отдача мала, она будет \textbf{понижать вероятности} соответствующих пар (состояние, действие).
  \end{itemize}
\end{itemize}

Для оценки градиента можно также \textbf{собирать несколько эпизодов (траекторий)}:

\CourseFigure{policy_gradient_multiple.png}{Обновление REINFORCE с усреднением по нескольким траекториям \protect\href{https://huggingface.co/datasets/huggingface-deep-rl-course/course-images/resolve/main/en/unit6/policy_gradient_multiple.png}{[оригинал]}}

\section{(Необязательно) Теорема о градиенте политики}\label{optional-the-policy-gradient-theorem}

В этом необязательном разделе мы \textbf{разберём, как дифференцировать целевую функцию для приближённого вычисления градиента политики}.

Для начала вспомним наши формулы:

\begin{enumerate}
\def\labelenumi{\arabic{enumi}.}
\tightlist
\item
  Целевая функция.
\end{enumerate}

\CourseFigure{expected_reward.png}{Ожидаемая отдача как целевая функция оптимизации \protect\href{https://huggingface.co/datasets/huggingface-deep-rl-course/course-images/resolve/main/en/unit6/expected_reward.png}{[оригинал]}}

\begin{enumerate}
\def\labelenumi{\arabic{enumi}.}
\setcounter{enumi}{1}
\tightlist
\item
  Вероятность траектории (при условии, что действия выбираются политикой \(\pi_\theta\)):
\end{enumerate}

\CourseFigure{probability.png}{Вероятность траектории при параметризованной политике \protect\href{https://huggingface.co/datasets/huggingface-deep-rl-course/course-images/resolve/main/en/unit6/probability.png}{[оригинал]}}

Итак, имеем:

\(\nabla_\theta J(\theta) =  \nabla_\theta \sum_{\tau}P(\tau;\theta)R(\tau)\)

Градиент суммы можно представить как сумму градиентов:

\(=  \sum_{\tau} \nabla_\theta (P(\tau;\theta)R(\tau)) = \sum_{\tau} \nabla_\theta P(\tau;\theta)R(\tau)\), поскольку \(R(\tau)\) не зависит от \(\theta\).

Затем умножим каждый член суммы на \(\frac{P(\tau;\theta)}{P(\tau;\theta)}\) (это допустимо, поскольку отношение равно 1).

\(= \sum_{\tau} \frac{P(\tau;\theta)}{P(\tau;\theta)}\nabla_\theta P(\tau;\theta)R(\tau)\)

Далее упростим выражение, так как \(\frac{P(\tau;\theta)}{P(\tau;\theta)}\nabla_\theta P(\tau;\theta) =  P(\tau;\theta)\frac{\nabla_\theta P(\tau;\theta)}{P(\tau;\theta)}\).

Поэтому сумму можно переписать следующим образом:

\(P(\tau;\theta)\frac{\nabla_\theta P(\tau;\theta)}{P(\tau;\theta)}= \sum_{\tau} P(\tau;\theta) \frac{\nabla_\theta P(\tau;\theta)}{P(\tau;\theta)}R(\tau)\)

Теперь воспользуемся \emph{трюком с производной логарифма} (derivative log trick), который также называют \emph{трюком отношения правдоподобий} (likelihood ratio trick) или \emph{трюком REINFORCE}. Это простое правило анализа: \(\nabla_x log f(x) = \frac{\nabla_x f(x)}{f(x)}\).

Поскольку в выражении присутствует \(\frac{\nabla_\theta P(\tau;\theta)}{P(\tau;\theta)}\), заменим его на \(\nabla_\theta log P(\tau|\theta)\).

Получаем градиент политики через правдоподобие:

\(\nabla_\theta J(\theta) = \sum_{\tau} P(\tau;\theta)  \nabla_\theta log P(\tau;\theta) R(\tau)\)

Благодаря новой формуле градиент можно оценивать по выборке траекторий (или, что то же самое, выборочно аппроксимировать градиент через отношение правдоподобий).

\(\nabla_\theta J(\theta) = \frac{1}{m} \sum^{m}_{i=1} \nabla_\theta log P(\tau^{(i)};\theta)R(\tau^{(i)})\), где каждая \(\tau^{(i)}\) --- выбранная из данных траектория.

Но математические преобразования ещё не закончены: требуется упростить \(\nabla_\theta log P(\tau|\theta)\).

Мы знаем, что:

\(\nabla_\theta log P(\tau^{(i)};\theta)= \nabla_\theta log[ \mu(s_0) \prod_{t=0}^{H} P(s_{t+1}^{(i)}|s_{t}^{(i)}, a_{t}^{(i)}) \pi_\theta(a_{t}^{(i)}|s_{t}^{(i)})]\)

Здесь \(\mu(s_0)\) --- распределение начальных состояний, а \(P(s_{t+1}^{(i)}|s_{t}^{(i)}, a_{t}^{(i)})\) --- динамика переходов MDP.

Логарифм произведения равен сумме логарифмов:

\(\nabla_\theta log P(\tau^{(i)};\theta)= \nabla_\theta \left[log \mu(s_0) + \sum\limits_{t=0}^{H}log P(s_{t+1}^{(i)}|s_{t}^{(i)} a_{t}^{(i)}) + \sum\limits_{t=0}^{H}log \pi_\theta(a_{t}^{(i)}|s_{t}^{(i)})\right]\)

Кроме того, градиент суммы равен сумме градиентов:

\(\nabla_\theta log P(\tau^{(i)};\theta)=\nabla_\theta log\mu(s_0) + \nabla_\theta \sum\limits_{t=0}^{H} log P(s_{t+1}^{(i)}|s_{t}^{(i)} a_{t}^{(i)}) + \nabla_\theta \sum\limits_{t=0}^{H} log \pi_\theta(a_{t}^{(i)}|s_{t}^{(i)})\)

Поскольку ни распределение начальных состояний, ни динамика переходов MDP не зависят от \(\theta\), производные обоих слагаемых равны нулю. Поэтому их можно убрать:

Поскольку \(\nabla_\theta \sum_{t=0}^{H} log P(s_{t+1}^{(i)}|s_{t}^{(i)}  a_{t}^{(i)}) = 0\) и \(\nabla_\theta \mu(s_0) = 0\).

\(\nabla_\theta log P(\tau^{(i)};\theta) =   \nabla_\theta \sum_{t=0}^{H} log \pi_\theta(a_{t}^{(i)}|s_{t}^{(i)})\)

Градиент суммы вновь перепишем как сумму градиентов:

\(\nabla_\theta log P(\tau^{(i)};\theta)=    \sum_{t=0}^{H} \nabla_\theta log \pi_\theta(a_{t}^{(i)}|s_{t}^{(i)})\)

Итак, окончательная формула оценки градиента политики:

\(\nabla_{\theta} J(\theta) = \hat{g} = \frac{1}{m} \sum^{m}_{i=1} \sum^{H}_{t=0} \nabla_\theta \log \pi_\theta(a^{(i)}_{t} | s_{t}^{(i)})R(\tau^{(i)})\)

\section{Практическое занятие}\label{hands-on}

\href{https://colab.research.google.com/github/huggingface/deep-rl-class/blob/main/notebooks/unit4/unit4.ipynb}{Официальный ноутбук модуля 4 (Google Colab)}

\href{http://hf.co/join/discord}{Помощь по курсу / Discord}

Теперь мы изучили теорию REINFORCE. \textbf{Пора реализовать агента на PyTorch} и проверить устойчивость его работы в CartPole-v1 и PixelCopter.

После этого ты сможешь самостоятельно доработать реализацию для более сложных сред.

\CourseFigure{envs.png}{Среды CartPole и PixelCopter (статичный кадр анимации) \protect\href{https://huggingface.co/datasets/huggingface-deep-rl-course/course-images/resolve/main/en/unit6/envs.gif}{[оригинал]}}

Для зачёта практического задания в системе сертификации необходимо загрузить обученные модели в Hub и:

\begin{itemize}
\tightlist
\item
  Получить результат \(\geq 350\) в \texttt{Cartpole-v1}.
\item
  Получить результат \(\geq 5\) в \texttt{PixelCopter}.
\end{itemize}

Чтобы узнать результат, открой таблицу лидеров и найди свою модель: \textbf{результат = mean\_reward - std награды}. \textbf{Если модель не отображается, перейди в конец страницы таблицы лидеров и нажми кнопку обновления}.

\textbf{Если модель не нашлась, прокрути страницу вниз и нажми кнопку обновления.}

Подробнее о сертификации читай в \href{https://huggingface.co/deep-rl-course/en/unit0/introduction#certification-process}{соответствующем разделе курса}.

Проверить свой прогресс можно здесь: \url{https://huggingface.co/spaces/ThomasSimonini/Check-my-progress-Deep-RL-Course}.

\textbf{Чтобы начать практику, нажми кнопку Open In Colab}:

\href{https://colab.research.google.com/github/huggingface/deep-rl-class/blob/master/notebooks/unit4/unit4.ipynb}{Открыть в Colab}

Мы настоятельно \textbf{рекомендуем студентам выполнять практические задания в Google Colab}, а не на личных компьютерах.

Используя Google Colab, \textbf{можно сосредоточиться на обучении и экспериментах, не отвлекаясь на технические детали} настройки окружений.

\subsection{Модуль 4: напиши первый алгоритм глубокого RL на PyTorch --- REINFORCE --- и проверь его устойчивость}\label{unit-4-code-your-first-deep-reinforcement-learning-algorithm-with-pytorch-reinforce.-and-test-its-robustness}

\CourseFigure{thumbnail.png}{Обложка модуля 4: градиент политики и REINFORCE \protect\href{https://huggingface.co/datasets/huggingface-deep-rl-course/course-images/resolve/main/en/unit6/thumbnail.png}{[оригинал]}}

В этом ноутбуке ты с нуля запрограммируешь свой первый алгоритм глубокого обучения с подкреплением --- REINFORCE (его также называют градиентом политики Монте-Карло).

REINFORCE относится к \emph{методам на основе политики}: он пытается \textbf{непосредственно оптимизировать политику без использования функции ценности действий}.

Точнее, REINFORCE --- это \emph{метод градиента политики}, подкласс методов на основе политики. Его цель --- \textbf{напрямую оптимизировать политику, оценивая веса оптимальной политики с помощью градиентного подъёма}.

Чтобы проверить устойчивость алгоритма, обучим его в двух разных простых средах:

\begin{itemize}
\tightlist
\item
  Cartpole-v1.
\item
  PixelcopterEnv.
\end{itemize}

Вот пример того, \textbf{что ты получишь к концу этого ноутбука.}

\CourseFigure{envs.png}{Среды CartPole и PixelCopter (статичный кадр анимации) \protect\href{https://huggingface.co/datasets/huggingface-deep-rl-course/course-images/resolve/main/en/unit6/envs.gif}{[оригинал]}}

\paragraph{Среды:}\label{environments}

\begin{itemize}
\tightlist
\item
  \href{https://www.gymlibrary.dev/environments/classic_control/cart_pole/}{CartPole-v1}
\item
  \href{https://pygame-learning-environment.readthedocs.io/en/latest/user/games/pixelcopter.html}{PixelCopter}
\end{itemize}

\paragraph{Библиотеки RL:}\label{rl-library}

\begin{itemize}
\tightlist
\item
  Python
\item
  PyTorch
\end{itemize}

Мы постоянно улучшаем наши учебные материалы. Поэтому \textbf{если заметишь проблемы в этом ноутбуке}, пожалуйста, \href{https://github.com/huggingface/deep-rl-class/issues}{создай issue в репозитории GitHub}.

\subsubsection{Цели ноутбука}\label{objectives-of-this-notebook}

После выполнения ноутбука ты сможешь:

\begin{itemize}
\tightlist
\item
  \textbf{Написать алгоритм REINFORCE с нуля на PyTorch.}
\item
  \textbf{Проверить устойчивость агента в простых средах.}
\item
  \textbf{Опубликовать обученного агента в Hugging Face Hub} вместе с видеозаписью и результатами оценки.
\end{itemize}

\subsubsection{Предварительные требования}\label{prerequisites}

Прежде чем приступать к ноутбуку, необходимо:

\href{https://huggingface.co/deep-rl-course/unit4/introduction}{Изучить градиенты политики в модуле 4}.

\subsection{Реализуем REINFORCE с нуля}\label{lets-code-reinforce-algorithm-from-scratch}

\subsubsection{Несколько советов}\label{some-advice-tip}

Лучше выполнять этот ноутбук в копии на своём Google Диске. Тогда, \textbf{если сеанс Colab завершится по тайм-ауту}, сохранённый ноутбук останется на Google Диске и не придётся начинать всё сначала.

Для этого нажми \texttt{Ctrl\ +\ S} или выбери \texttt{File\ \textgreater{}\ Save\ a\ copy\ in\ Google\ Drive}.

\subsubsection{Подключаем GPU}\label{set-the-gpu}

\begin{itemize}
\tightlist
\item
  Чтобы \textbf{ускорить обучение агента, используем GPU}. Для этого выбери \texttt{Runtime\ \textgreater{}\ Change\ Runtime\ type}.
\end{itemize}

\CourseFigure{gpu-step1.jpg}{Выбор режима выполнения Colab: первый шаг \protect\href{https://huggingface.co/datasets/huggingface-deep-rl-course/course-images/resolve/main/en/notebooks/gpu-step1.jpg}{[оригинал]}}

\begin{itemize}
\tightlist
\item
  \texttt{Hardware\ Accelerator\ \textgreater{}\ GPU}.
\end{itemize}

\CourseFigure{gpu-step2.jpg}{Выбор аппаратного ускорителя GPU в Colab \protect\href{https://huggingface.co/datasets/huggingface-deep-rl-course/course-images/resolve/main/en/notebooks/gpu-step2.jpg}{[оригинал]}}

\subsubsection{Создаём виртуальный дисплей}\label{create-a-virtual-display}

Во время работы с ноутбуком нам понадобится создавать видеозаписи прохождений. Для этого в Colab \textbf{нужен виртуальный экран, позволяющий отображать среду} (и сохранять её кадры).

Следующая ячейка установит библиотеки, создаст виртуальный экран и запустит его.

\begin{lstlisting}
%%capture
!apt install python-opengl
!apt install ffmpeg
!apt install xvfb
!pip install pyvirtualdisplay
!pip install pyglet==1.5.1
\end{lstlisting}

\begin{lstlisting}
# Virtual display
from pyvirtualdisplay import Display

virtual_display = Display(visible=0, size=(1400, 900))
virtual_display.start()
\end{lstlisting}

\subsubsection{Устанавливаем зависимости}\label{install-the-dependencies}

Сначала установим необходимые зависимости. Их несколько:

\begin{itemize}
\tightlist
\item
  \texttt{gym}.
\item
  \texttt{gym-games}: дополнительные среды Gym, созданные с помощью PyGame.
\item
  \texttt{huggingface\_hub}: Hub --- центральная площадка, на которой можно публиковать и изучать модели и датасеты. Здесь есть версионирование, метрики, визуализации и другие возможности для совместной работы.
\end{itemize}

Возможно, ты удивляешься, почему мы устанавливаем Gym, а не более новую Gymnasium. \textbf{Потому что используемая нами библиотека gym-games ещё не была адаптирована для Gymnasium}.

Основные различия, с которыми ты столкнёшься:

\begin{itemize}
\tightlist
\item
  В \texttt{gym} вместо отдельных флагов \texttt{terminated} и \texttt{truncated} используется единственный флаг \texttt{done}.
\item
  В \texttt{gym} вызов \texttt{env.step()} возвращает \texttt{state,\ reward,\ done,\ info}.
\end{itemize}

Подробнее о различиях Gym и Gymnasium можно узнать здесь: \url{https://gymnasium.farama.org/content/migration-guide/}.

Все доступные модели REINFORCE можно посмотреть здесь: \url{https://huggingface.co/models?other=reinforce}.

А все модели глубокого обучения с подкреплением --- здесь: \url{https://huggingface.co/models?pipeline_tag=reinforcement-learning}.

\begin{lstlisting}
!pip install -r https://raw.githubusercontent.com/huggingface/deep-rl-class/main/notebooks/unit4/requirements-unit4.txt
\end{lstlisting}

\subsubsection{Импортируем пакеты}\label{import-the-packages}

Помимо установленных библиотек, мы импортируем:

\begin{itemize}
\tightlist
\item
  \texttt{imageio}: библиотеку, которая поможет создать видеозапись прохождения.
\end{itemize}

\begin{lstlisting}
import numpy as np

from collections import deque

import matplotlib.pyplot as plt
%matplotlib inline

# PyTorch
import torch
import torch.nn as nn
import torch.nn.functional as F
import torch.optim as optim
from torch.distributions import Categorical

# Gym
import gym
import gym_pygame

# Hugging Face Hub
from huggingface_hub import notebook_login # To log to our Hugging Face account to be able to upload models to the Hub.
import imageio
\end{lstlisting}

\subsubsection{Проверяем наличие GPU}\label{check-if-we-have-a-gpu}

\begin{itemize}
\tightlist
\item
  Проверим, доступен ли GPU.
\item
  Если да, ты увидишь \texttt{device:cuda0}.
\end{itemize}

\begin{lstlisting}
device = torch.device("cuda:0" if torch.cuda.is_available() else "cpu")
\end{lstlisting}

\begin{lstlisting}
print(device)
\end{lstlisting}

Теперь всё готово для реализации алгоритма REINFORCE.

\subsection{Первый агент: играем в CartPole-v1}\label{first-agent-playing-cartpole-v1}

\subsubsection{Создаём среду CartPole и разбираемся в её работе}\label{create-the-cartpole-environment-and-understand-how-it-works}

\paragraph{\texorpdfstring{\href{https://www.gymlibrary.dev/environments/classic_control/cart_pole/}{Среда}}{Среда}}\label{the-environment}

\paragraph{Почему мы используем простую среду вроде CartPole-v1?}\label{why-do-we-use-a-simple-environment-like-cartpole-v1}

Как объясняется в \href{https://stable-baselines3.readthedocs.io/en/master/guide/rl_tips.html}{советах по RL}, когда реализуешь агента с нуля, нужно \textbf{сначала убедиться, что он работает правильно, и отладить его в простых средах}. Так намного легче искать ошибки, прежде чем переходить к сложным задачам.

\begin{quote}
Попробуй получить хотя бы какие-то «признаки жизни» на игрушечных задачах.
\end{quote}

\begin{quote}
Проверяй реализацию, запуская её во всё более сложных средах (результаты можно сравнивать с RL Zoo). На этом этапе обычно приходится подбирать гиперпараметры.
\end{quote}

\paragraph{Среда CartPole-v1}\label{the-cartpole-v1-environment}

\begin{quote}
К тележке, движущейся по рельсу без трения, через неуправляемый шарнир прикреплён шест. В начале шест расположен вертикально. Задача --- удерживать его в равновесии, прикладывая к тележке силы влево или вправо.
\end{quote}

Итак, начинаем с CartPole-v1. Нужно перемещать тележку влево или вправо \textbf{так, чтобы шест не падал и оставался в равновесии}.

Эпизод заканчивается, если:

\begin{itemize}
\tightlist
\item
  Угол отклонения шеста превышает \(\pm12^{\circ}\).
\item
  Координата тележки выходит за пределы \(\pm2{,}4\).
\item
  Длительность эпизода превышает 500 шагов.
\end{itemize}

На каждом шаге, пока шест удерживается в равновесии, агент получает награду +1.

\begin{lstlisting}
env_id = "CartPole-v1"
# Create the env
env = gym.make(env_id)

# Create the evaluation env
eval_env = gym.make(env_id)

# Get the state space and action space
s_size = env.observation_space.shape[0]
a_size = env.action_space.n
\end{lstlisting}

\begin{lstlisting}
print("_____OBSERVATION SPACE_____ \n")
print("The State Space is: ", s_size)
print("Sample observation", env.observation_space.sample())  # Get a random observation
\end{lstlisting}

\begin{lstlisting}
print("\n _____ACTION SPACE_____ \n")
print("The Action Space is: ", a_size)
print("Action Space Sample", env.action_space.sample())  # Take a random action
\end{lstlisting}

\subsubsection{Построим архитектуру REINFORCE}\label{lets-build-the-reinforce-architecture}

Эта реализация основана на трёх примерах:

\begin{itemize}
\tightlist
\item
  \href{https://github.com/pytorch/examples/blob/main/reinforcement_learning/reinforce.py}{Официальный пример RL от PyTorch}.
\item
  \href{https://github.com/udacity/deep-reinforcement-learning/blob/master/reinforce/REINFORCE.ipynb}{REINFORCE от Udacity}.
\item
  \href{https://github.com/huggingface/deep-rl-class/pull/95}{Улучшение интеграции от Chris1nexus}.
\end{itemize}

\CourseFigure{reinforce.png}{Архитектура нейросети политики REINFORCE \protect\href{https://huggingface.co/datasets/huggingface-deep-rl-course/course-images/resolve/main/en/unit6/reinforce.png}{[оригинал]}}

Итак, нам нужно:

\begin{itemize}
\tightlist
\item
  Два полносвязных слоя (fc1 и fc2).
\item
  Использовать ReLU в качестве функции активации слоя fc1.
\item
  Использовать Softmax, чтобы получить распределение вероятностей действий на выходе.
\end{itemize}

\begin{lstlisting}
class Policy(nn.Module):
    def __init__(self, s_size, a_size, h_size):
        super(Policy, self).__init__()
        # Create two fully connected layers



    def forward(self, x):
        # Define the forward pass
        # state goes to fc1 then we apply ReLU activation function

        # fc1 outputs goes to fc2

        # We output the softmax

    def act(self, state):
        """
        Given a state, take action
        """
        state = torch.from_numpy(state).float().unsqueeze(0).to(device)
        probs = self.forward(state).cpu()
        m = Categorical(probs)
        action = np.argmax(m)
        return action.item(), m.log_prob(action)
\end{lstlisting}

\paragraph{Решение}\label{solution}

\begin{lstlisting}
class Policy(nn.Module):
    def __init__(self, s_size, a_size, h_size):
        super(Policy, self).__init__()
        self.fc1 = nn.Linear(s_size, h_size)
        self.fc2 = nn.Linear(h_size, a_size)

    def forward(self, x):
        x = F.relu(self.fc1(x))
        x = self.fc2(x)
        return F.softmax(x, dim=1)

    def act(self, state):
        state = torch.from_numpy(state).float().unsqueeze(0).to(device)
        probs = self.forward(state).cpu()
        m = Categorical(probs)
        action = np.argmax(m)
        return action.item(), m.log_prob(action)
\end{lstlisting}

Я допустил ошибку. Сможешь догадаться, где именно?

\begin{itemize}
\tightlist
\item
  Чтобы выяснить это, выполним прямой проход сети:
\end{itemize}

\begin{lstlisting}
debug_policy = Policy(s_size, a_size, 64).to(device)
debug_policy.act(env.reset())
\end{lstlisting}

\begin{itemize}
\item
  Как видно, появляется ошибка \texttt{ValueError:\ The\ value\ argument\ to\ log\_prob\ must\ be\ a\ Tensor}.
\item
  Она означает, что \texttt{action} при вызове \texttt{m.log\_prob(action)} должен быть тензором (Tensor), \textbf{но это не так}.
\item
  Почему так происходит? Посмотри на функцию act и попробуй понять, почему она не работает.
\end{itemize}

Совет: в этой реализации есть ошибка. Помни, что функция act должна \textbf{выбирать случайное действие из распределения вероятностей действий}.

\paragraph{Настоящее решение}\label{real-solution}

\begin{lstlisting}
class Policy(nn.Module):
    def __init__(self, s_size, a_size, h_size):
        super(Policy, self).__init__()
        self.fc1 = nn.Linear(s_size, h_size)
        self.fc2 = nn.Linear(h_size, a_size)

    def forward(self, x):
        x = F.relu(self.fc1(x))
        x = self.fc2(x)
        return F.softmax(x, dim=1)

    def act(self, state):
        state = torch.from_numpy(state).float().unsqueeze(0).to(device)
        probs = self.forward(state).cpu()
        m = Categorical(probs)
        action = m.sample()
        return action.item(), m.log_prob(action)
\end{lstlisting}

Использование CartPole облегчило отладку, поскольку \textbf{мы знаем, что ошибка вызвана нашей реализацией, а не простой средой}.

\begin{itemize}
\item
  Поскольку \textbf{мы хотим выбирать действие из распределения вероятностей}, нельзя использовать \texttt{action\ =\ np.argmax(m)}: этот вызов всегда возвращает действие с наибольшей вероятностью.
\item
  Нужно заменить его на \texttt{action\ =\ m.sample()}, который выбирает действие случайным образом согласно распределению \(P(\cdot\mid s)\).
\end{itemize}

\paragraph{Реализуем алгоритм обучения REINFORCE}\label{lets-build-the-reinforce-training-algorithm}

Псевдокод алгоритма REINFORCE:

\CourseFigure{pg_pseudocode.png}{Псевдокод алгоритма REINFORCE \protect\href{https://huggingface.co/datasets/huggingface-deep-rl-course/course-images/resolve/main/en/unit6/pg_pseudocode.png}{[оригинал]}}

\begin{itemize}
\item
  При вычислении отдачи \(G_t\) (строка 6) видно, что мы суммируем дисконтированные награды, \textbf{начиная с момента времени \(t\)}.
\item
  Почему? Потому что политика должна \textbf{усиливать действия на основании их последствий}: награды, полученные до действия, не имеют значения (они не вызваны этим действием), \textbf{важны только последующие награды}.
\item
  Перед реализацией прочитай раздел \href{https://spinningup.openai.com/en/latest/spinningup/rl_intro3.html\#don-t-let-the-past-distract-you}{«Не позволяй прошлому отвлекать тебя»}, где объясняется использование градиента политики на основе будущих наград (reward-to-go).
\end{itemize}

Мы используем интересный приём, предложенный \href{https://github.com/Chris1nexus}{Chris1nexus}, чтобы \textbf{эффективно вычислять отдачу на каждом временном шаге}. Процедура объяснена в комментариях к коду. Также можно \href{https://github.com/huggingface/deep-rl-class/pull/95}{прочитать пояснения к PR}. Главная идея --- \textbf{вычислять отдачу для каждого шага без лишних затрат}.

Второй вопрос, который может возникнуть: \textbf{почему мы минимизируем функцию потерь}? Разве ранее мы не говорили о градиентном подъёме, а не о спуске?

\begin{itemize}
\tightlist
\item
  Нам нужно максимизировать функцию полезности \(J(\theta)\), но в PyTorch и TensorFlow удобнее \textbf{минимизировать целевую функцию}.

  \begin{itemize}
  \tightlist
  \item
    Допустим, на некотором шаге мы хотим усилить действие 3. До обучения его вероятность равна 0,25.
  \item
    Нужно изменить \(\theta\) так, чтобы \(\pi_\theta(a_3|s; \theta) > 0.25\).
  \item
    Поскольку сумма всех вероятностей равна 1, максимизация \(\pi_\theta(a_3|s; \theta)\) будет \textbf{уменьшать вероятности других действий}.
  \item
    Поэтому нужно указать PyTorch \textbf{минимизировать \(1 - \pi_\theta(a_3|s; \theta)\)}.
  \item
    Такая функция потерь стремится к нулю при приближении \(\pi_\theta(a_3|s; \theta)\) к 1.
  \item
    Тем самым мы направляем градиент на увеличение \(\pi_\theta(a_3|s; \theta)\).
  \end{itemize}
\end{itemize}

\begin{lstlisting}
def reinforce(policy, optimizer, n_training_episodes, max_t, gamma, print_every):
    # Help us to calculate the score during the training
    scores_deque = deque(maxlen=100)
    scores = []
    # Line 3 of pseudocode
    for i_episode in range(1, n_training_episodes+1):
        saved_log_probs = []
        rewards = []
        state = # TODO: reset the environment
        # Line 4 of pseudocode
        for t in range(max_t):
            action, log_prob = # TODO get the action
            saved_log_probs.append(log_prob)
            state, reward, done, _ = # TODO: take an env step
            rewards.append(reward)
            if done:
                break
        scores_deque.append(sum(rewards))
        scores.append(sum(rewards))

        # Line 6 of pseudocode: calculate the return
        returns = deque(maxlen=max_t)
        n_steps = len(rewards)
        # Compute the discounted returns at each timestep,
        # as the sum of the gamma-discounted return at time t (G_t) + the reward at time t

        # In O(N) time, where N is the number of time steps
        # (this definition of the discounted return G_t follows the definition of this quantity
        # shown at page 44 of Sutton&Barto 2017 2nd draft)
        # G_t = r_(t+1) + r_(t+2) + ...

        # Given this formulation, the returns at each timestep t can be computed
        # by re-using the computed future returns G_(t+1) to compute the current return G_t
        # G_t = r_(t+1) + gamma*G_(t+1)
        # G_(t-1) = r_t + gamma* G_t
        # (this follows a dynamic programming approach, with which we memorize solutions in order
        # to avoid computing them multiple times)

        # This is correct since the above is equivalent to (see also page 46 of Sutton&Barto 2017 2nd draft)
        # G_(t-1) = r_t + gamma*r_(t+1) + gamma*gamma*r_(t+2) + ...


        ## Given the above, we calculate the returns at timestep t as:
        #               gamma[t] * return[t] + reward[t]
        #
        ## We compute this starting from the last timestep to the first, in order
        ## to employ the formula presented above and avoid redundant computations that would be needed
        ## if we were to do it from first to last.

        ## Hence, the queue "returns" will hold the returns in chronological order, from t=0 to t=n_steps
        ## thanks to the appendleft() function which allows to append to the position 0 in constant time O(1)
        ## a normal python list would instead require O(N) to do this.
        for t in range(n_steps)[::-1]:
            disc_return_t = (returns[0] if len(returns)>0 else 0)
            returns.appendleft(    ) # TODO: complete here

        ## standardization of the returns is employed to make training more stable
        eps = np.finfo(np.float32).eps.item()

        ## eps is the smallest representable float, which is
        # added to the standard deviation of the returns to avoid numerical instabilities
        returns = torch.tensor(returns)
        returns = (returns - returns.mean()) / (returns.std() + eps)

        # Line 7:
        policy_loss = []
        for log_prob, disc_return in zip(saved_log_probs, returns):
            policy_loss.append(-log_prob * disc_return)
        policy_loss = torch.cat(policy_loss).sum()

        # Line 8: PyTorch prefers gradient descent
        optimizer.zero_grad()
        policy_loss.backward()
        optimizer.step()

        if i_episode % print_every == 0:
            print('Episode {}\tAverage Score: {:.2f}'.format(i_episode, np.mean(scores_deque)))

    return scores
\end{lstlisting}

\subparagraph{Решение}\label{solution-1}

\begin{lstlisting}
def reinforce(policy, optimizer, n_training_episodes, max_t, gamma, print_every):
    # Help us to calculate the score during the training
    scores_deque = deque(maxlen=100)
    scores = []
    # Line 3 of pseudocode
    for i_episode in range(1, n_training_episodes + 1):
        saved_log_probs = []
        rewards = []
        state = env.reset()
        # Line 4 of pseudocode
        for t in range(max_t):
            action, log_prob = policy.act(state)
            saved_log_probs.append(log_prob)
            state, reward, done, _ = env.step(action)
            rewards.append(reward)
            if done:
                break
        scores_deque.append(sum(rewards))
        scores.append(sum(rewards))

        # Line 6 of pseudocode: calculate the return
        returns = deque(maxlen=max_t)
        n_steps = len(rewards)
        # Compute the discounted returns at each timestep,
        # as
        #      the sum of the gamma-discounted return at time t (G_t) + the reward at time t
        #
        # In O(N) time, where N is the number of time steps
        # (this definition of the discounted return G_t follows the definition of this quantity
        # shown at page 44 of Sutton&Barto 2017 2nd draft)
        # G_t = r_(t+1) + r_(t+2) + ...

        # Given this formulation, the returns at each timestep t can be computed
        # by re-using the computed future returns G_(t+1) to compute the current return G_t
        # G_t = r_(t+1) + gamma*G_(t+1)
        # G_(t-1) = r_t + gamma* G_t
        # (this follows a dynamic programming approach, with which we memorize solutions in order
        # to avoid computing them multiple times)

        # This is correct since the above is equivalent to (see also page 46 of Sutton&Barto 2017 2nd draft)
        # G_(t-1) = r_t + gamma*r_(t+1) + gamma*gamma*r_(t+2) + ...

        ## Given the above, we calculate the returns at timestep t as:
        #               gamma[t] * return[t] + reward[t]
        #
        ## We compute this starting from the last timestep to the first, in order
        ## to employ the formula presented above and avoid redundant computations that would be needed
        ## if we were to do it from first to last.

        ## Hence, the queue "returns" will hold the returns in chronological order, from t=0 to t=n_steps
        ## thanks to the appendleft() function which allows to append to the position 0 in constant time O(1)
        ## a normal python list would instead require O(N) to do this.
        for t in range(n_steps)[::-1]:
            disc_return_t = returns[0] if len(returns) > 0 else 0
            returns.appendleft(gamma * disc_return_t + rewards[t])

        ## standardization of the returns is employed to make training more stable
        eps = np.finfo(np.float32).eps.item()
        ## eps is the smallest representable float, which is
        # added to the standard deviation of the returns to avoid numerical instabilities
        returns = torch.tensor(returns)
        returns = (returns - returns.mean()) / (returns.std() + eps)

        # Line 7:
        policy_loss = []
        for log_prob, disc_return in zip(saved_log_probs, returns):
            policy_loss.append(-log_prob * disc_return)
        policy_loss = torch.cat(policy_loss).sum()

        # Line 8: PyTorch prefers gradient descent
        optimizer.zero_grad()
        policy_loss.backward()
        optimizer.step()

        if i_episode % print_every == 0:
            print("Episode {}\tAverage Score: {:.2f}".format(i_episode, np.mean(scores_deque)))

    return scores
\end{lstlisting}

\subsubsection{Обучаем агента}\label{train-it}

\begin{itemize}
\tightlist
\item
  Теперь всё готово к обучению агента.
\item
  Но сначала зададим переменную, содержащую все гиперпараметры обучения.
\item
  Ты можешь (и даже должен) изменять параметры обучения.
\end{itemize}

\begin{lstlisting}
cartpole_hyperparameters = {
    "h_size": 16,
    "n_training_episodes": 1000,
    "n_evaluation_episodes": 10,
    "max_t": 1000,
    "gamma": 1.0,
    "lr": 1e-2,
    "env_id": env_id,
    "state_space": s_size,
    "action_space": a_size,
}
\end{lstlisting}

\begin{lstlisting}
# Create policy and place it to the device
cartpole_policy = Policy(
    cartpole_hyperparameters["state_space"],
    cartpole_hyperparameters["action_space"],
    cartpole_hyperparameters["h_size"],
).to(device)
cartpole_optimizer = optim.Adam(cartpole_policy.parameters(), lr=cartpole_hyperparameters["lr"])
\end{lstlisting}

\begin{lstlisting}
scores = reinforce(
    cartpole_policy,
    cartpole_optimizer,
    cartpole_hyperparameters["n_training_episodes"],
    cartpole_hyperparameters["max_t"],
    cartpole_hyperparameters["gamma"],
    100,
)
\end{lstlisting}

\subsubsection{Определяем метод оценки}\label{define-evaluation-method}

\begin{itemize}
\tightlist
\item
  Здесь мы определим метод оценки, с помощью которого будем тестировать нашего агента REINFORCE.
\end{itemize}

\begin{lstlisting}
def evaluate_agent(env, max_steps, n_eval_episodes, policy):
    """
    Evaluate the agent for ``n_eval_episodes`` episodes and returns average reward and std of reward.
    :param env: The evaluation environment
    :param n_eval_episodes: Number of episode to evaluate the agent
    :param policy: The Reinforce agent
    """
    episode_rewards = []
    for episode in range(n_eval_episodes):
        state = env.reset()
        step = 0
        done = False
        total_rewards_ep = 0

        for step in range(max_steps):
            action, _ = policy.act(state)
            new_state, reward, done, info = env.step(action)
            total_rewards_ep += reward

            if done:
                break
            state = new_state
        episode_rewards.append(total_rewards_ep)
    mean_reward = np.mean(episode_rewards)
    std_reward = np.std(episode_rewards)

    return mean_reward, std_reward
\end{lstlisting}

\subsubsection{Оцениваем работу агента}\label{evaluate-our-agent}

\begin{lstlisting}
evaluate_agent(
    eval_env, cartpole_hyperparameters["max_t"], cartpole_hyperparameters["n_evaluation_episodes"], cartpole_policy
)
\end{lstlisting}

\paragraph{Публикуем обученную модель на Hub}\label{publish-our-trained-model-on-the-hub}

Теперь, убедившись, что обучение дало хорошие результаты, мы можем опубликовать обученную модель на Hub одной строкой кода.

Вот пример карточки модели (Model Card):

\CourseFigure{modelcard.png}{Интерфейс карточки модели в Hugging Face \protect\href{https://huggingface.co/datasets/huggingface-deep-rl-course/course-images/resolve/main/en/unit6/modelcard.png}{[оригинал]}}

\paragraph{Загружаем модель на Hub}\label{push-to-the-hub}

\subparagraph{Не изменяй этот код}\label{do-not-modify-this-code}

\begin{lstlisting}
from huggingface_hub import HfApi, snapshot_download
from huggingface_hub.repocard import metadata_eval_result, metadata_save

from pathlib import Path
import datetime
import json
import imageio

import tempfile

import os
\end{lstlisting}

\begin{lstlisting}
def record_video(env, policy, out_directory, fps=30):
    """
    Generate a replay video of the agent
    :param env
    :param Qtable: Qtable of our agent
    :param out_directory
    :param fps: how many frame per seconds (with taxi-v3 and frozenlake-v1 we use 1)
    """
    images = []
    done = False
    state = env.reset()
    img = env.render(mode="rgb_array")
    images.append(img)
    while not done:
        # Take the action (index) that have the maximum expected future reward given that state
        action, _ = policy.act(state)
        state, reward, done, info = env.step(action)  # We directly put next_state = state for recording logic
        img = env.render(mode="rgb_array")
        images.append(img)
    imageio.mimsave(out_directory, [np.array(img) for i, img in enumerate(images)], fps=fps)
\end{lstlisting}

\begin{lstlisting}
def push_to_hub(repo_id,
                model,
                hyperparameters,
                eval_env,
                video_fps=30
                ):
  """
  Evaluate, Generate a video and Upload a model to Hugging Face Hub.
  This method does the complete pipeline:
  - It evaluates the model
  - It generates the model card
  - It generates a replay video of the agent
  - It pushes everything to the Hub

  :param repo_id: repo_id: id of the model repository from the Hugging Face Hub
  :param model: the pytorch model we want to save
  :param hyperparameters: training hyperparameters
  :param eval_env: evaluation environment
  :param video_fps: how many frame per seconds to record our video replay
  """

  _, repo_name = repo_id.split("/")
  api = HfApi()

  # Step 1: Create the repo
  repo_url = api.create_repo(
        repo_id=repo_id,
        exist_ok=True,
  )

  with tempfile.TemporaryDirectory() as tmpdirname:
    local_directory = Path(tmpdirname)

    # Step 2: Save the model
    torch.save(model, local_directory / "model.pt")

    # Step 3: Save the hyperparameters to JSON
    with open(local_directory / "hyperparameters.json", "w") as outfile:
      json.dump(hyperparameters, outfile)

    # Step 4: Evaluate the model and build JSON
    mean_reward, std_reward = evaluate_agent(eval_env,
                                            hyperparameters["max_t"],
                                            hyperparameters["n_evaluation_episodes"],
                                            model)
    # Get datetime
    eval_datetime = datetime.datetime.now()
    eval_form_datetime = eval_datetime.isoformat()

    evaluate_data = {
          "env_id": hyperparameters["env_id"],
          "mean_reward": mean_reward,
          "n_evaluation_episodes": hyperparameters["n_evaluation_episodes"],
          "eval_datetime": eval_form_datetime,
    }

    # Write a JSON file
    with open(local_directory / "results.json", "w") as outfile:
        json.dump(evaluate_data, outfile)

    # Step 5: Create the model card
    env_name = hyperparameters["env_id"]

    metadata = {}
    metadata["tags"] = [
          env_name,
          "reinforce",
          "reinforcement-learning",
          "custom-implementation",
          "deep-rl-class"
      ]

    # Add metrics
    eval = metadata_eval_result(
        model_pretty_name=repo_name,
        task_pretty_name="reinforcement-learning",
        task_id="reinforcement-learning",
        metrics_pretty_name="mean_reward",
        metrics_id="mean_reward",
        metrics_value=f"{mean_reward:.2f} +/- {std_reward:.2f}",
        dataset_pretty_name=env_name,
        dataset_id=env_name,
      )

    # Merges both dictionaries
    metadata = {**metadata, **eval}

    model_card = f"""
  # **Reinforce** Agent playing **{env_id}**
  This is a trained model of a **Reinforce** agent playing **{env_id}** .
  To learn to use this model and train yours check Unit 4 of the Deep Reinforcement Learning Course: https://huggingface.co/deep-rl-course/unit4/introduction
  """

    readme_path = local_directory / "README.md"
    readme = ""
    if readme_path.exists():
        with readme_path.open("r", encoding="utf8") as f:
          readme = f.read()
    else:
      readme = model_card

    with readme_path.open("w", encoding="utf-8") as f:
      f.write(readme)

    # Save our metrics to Readme metadata
    metadata_save(readme_path, metadata)

    # Step 6: Record a video
    video_path =  local_directory / "replay.mp4"
    record_video(env, model, video_path, video_fps)

    # Step 7. Push everything to the Hub
    api.upload_folder(
          repo_id=repo_id,
          folder_path=local_directory,
          path_in_repo=".",
    )

    print(f"Your model is pushed to the Hub. You can view your model here: {repo_url}")
\end{lstlisting}

С помощью \texttt{push\_to\_hub} \textbf{ты можешь оценить агента, записать видео его работы, создать карточку модели и загрузить её на Hub}.

Таким образом:

\begin{itemize}
\tightlist
\item
  Ты сможешь \textbf{показать результаты своей работы}.
\item
  Ты сможешь \textbf{посмотреть, как играет твой агент}.
\item
  Ты сможешь \textbf{поделиться агентом с сообществом, чтобы другие тоже могли его использовать}.
\item
  Ты сможешь \textbf{открыть таблицу лидеров и сравнить качество работы своего агента с результатами однокурсников}: \url{https://huggingface.co/spaces/huggingface-projects/Deep-Reinforcement-Learning-Leaderboard}.
\end{itemize}

Чтобы поделиться моделью с сообществом, осталось выполнить ещё три шага:

\begin{enumerate}
\def\labelenumi{\arabic{enumi}.}
\item
  Если учётной записи Hugging Face пока нет, создай её: \url{https://huggingface.co/join}.
\item
  Войди в аккаунт, затем сохрани токен аутентификации с сайта Hugging Face.
\end{enumerate}

\begin{itemize}
\tightlist
\item
  Создай новый токен на странице \url{https://huggingface.co/settings/tokens} \textbf{с правами записи (write)}.
\end{itemize}

\CourseFigure{create-token.jpg}{Создание токена доступа Hugging Face \protect\href{https://huggingface.co/datasets/huggingface-deep-rl-course/course-images/resolve/main/en/notebooks/create-token.jpg}{[оригинал]}}

\begin{lstlisting}
notebook_login()
\end{lstlisting}

Если ты не используешь Google Colab или Jupyter Notebook, вместо этого выполни команду \texttt{huggingface-cli\ login} (или \texttt{login}).

\begin{enumerate}
\def\labelenumi{\arabic{enumi}.}
\setcounter{enumi}{2}
\tightlist
\item
  Теперь мы готовы отправить обученного агента в Hub. Для этого используем функцию \texttt{package\_to\_hub()}.
\end{enumerate}

\begin{lstlisting}
repo_id = ""  # TODO Define your repo id {username/Reinforce-{model-id}}
push_to_hub(
    repo_id,
    cartpole_policy,  # The model we want to save
    cartpole_hyperparameters,  # Hyperparameters
    eval_env,  # Evaluation environment
    video_fps=30
)
\end{lstlisting}

Проверив устойчивость нашей реализации, перейдём к более сложной среде --- PixelCopter.

\subsubsection{Второй агент: PixelCopter}\label{second-agent-pixelcopter}

\paragraph{Изучаем среду PixelCopter}\label{study-the-pixelcopter-environment}

\begin{itemize}
\tightlist
\item
  \href{https://pygame-learning-environment.readthedocs.io/en/latest/user/games/pixelcopter.html}{Документация среды}.
\end{itemize}

\begin{lstlisting}
env_id = "Pixelcopter-PLE-v0"
env = gym.make(env_id)
eval_env = gym.make(env_id)
s_size = env.observation_space.shape[0]
a_size = env.action_space.n
\end{lstlisting}

\begin{lstlisting}
print("_____OBSERVATION SPACE_____ \n")
print("The State Space is: ", s_size)
print("Sample observation", env.observation_space.sample())  # Get a random observation
\end{lstlisting}

\begin{lstlisting}
print("\n _____ACTION SPACE_____ \n")
print("The Action Space is: ", a_size)
print("Action Space Sample", env.action_space.sample())  # Take a random action
\end{lstlisting}

Пространство наблюдений (7 компонентов):

\begin{itemize}
\tightlist
\item
  Координата игрока по оси \(y\).
\item
  Скорость игрока.
\item
  Расстояние от игрока до пола.
\item
  Расстояние от игрока до потолка.
\item
  Расстояние по оси \(x\) от игрока до следующего препятствия.
\item
  Координата верхней границы следующего препятствия по оси \(y\).
\item
  Координата нижней границы следующего препятствия по оси \(y\).
\end{itemize}

Пространство действий (2 действия):

\begin{itemize}
\tightlist
\item
  Подняться вверх (нажать ускоритель).
\item
  Ничего не делать (не нажимать ускоритель).
\end{itemize}

Функция вознаграждения:

\begin{itemize}
\tightlist
\item
  За прохождение каждого вертикального препятствия агент получает положительную награду +1. При достижении терминального состояния он получает отрицательную награду -1.
\end{itemize}

\paragraph{Определяем новую политику}\label{define-the-new-policy}

\begin{itemize}
\tightlist
\item
  Нам нужна более глубокая нейронная сеть, поскольку среда стала сложнее.
\end{itemize}

\begin{lstlisting}
class Policy(nn.Module):
    def __init__(self, s_size, a_size, h_size):
        super(Policy, self).__init__()
        # Define the three layers here

    def forward(self, x):
        # Define the forward process here
        return F.softmax(x, dim=1)

    def act(self, state):
        state = torch.from_numpy(state).float().unsqueeze(0).to(device)
        probs = self.forward(state).cpu()
        m = Categorical(probs)
        action = m.sample()
        return action.item(), m.log_prob(action)
\end{lstlisting}

\subparagraph{Решение}\label{solution-2}

\begin{lstlisting}
class Policy(nn.Module):
    def __init__(self, s_size, a_size, h_size):
        super(Policy, self).__init__()
        self.fc1 = nn.Linear(s_size, h_size)
        self.fc2 = nn.Linear(h_size, h_size * 2)
        self.fc3 = nn.Linear(h_size * 2, a_size)

    def forward(self, x):
        x = F.relu(self.fc1(x))
        x = F.relu(self.fc2(x))
        x = self.fc3(x)
        return F.softmax(x, dim=1)

    def act(self, state):
        state = torch.from_numpy(state).float().unsqueeze(0).to(device)
        probs = self.forward(state).cpu()
        m = Categorical(probs)
        action = m.sample()
        return action.item(), m.log_prob(action)
\end{lstlisting}

\paragraph{Задаём гиперпараметры}\label{define-the-hyperparameters}

\begin{itemize}
\tightlist
\item
  Потому что эта среда сложнее.
\item
  Особенно важно увеличить число нейронов в скрытых слоях.
\end{itemize}

\begin{lstlisting}
pixelcopter_hyperparameters = {
    "h_size": 64,
    "n_training_episodes": 50000,
    "n_evaluation_episodes": 10,
    "max_t": 10000,
    "gamma": 0.99,
    "lr": 1e-4,
    "env_id": env_id,
    "state_space": s_size,
    "action_space": a_size,
}
\end{lstlisting}

\paragraph{Обучаем агента}\label{train-it-1}

\begin{itemize}
\tightlist
\item
  Теперь всё готово для обучения агента.
\end{itemize}

\begin{lstlisting}
# Create policy and place it to the device
# torch.manual_seed(50)
pixelcopter_policy = Policy(
    pixelcopter_hyperparameters["state_space"],
    pixelcopter_hyperparameters["action_space"],
    pixelcopter_hyperparameters["h_size"],
).to(device)
pixelcopter_optimizer = optim.Adam(pixelcopter_policy.parameters(), lr=pixelcopter_hyperparameters["lr"])
\end{lstlisting}

\begin{lstlisting}
scores = reinforce(
    pixelcopter_policy,
    pixelcopter_optimizer,
    pixelcopter_hyperparameters["n_training_episodes"],
    pixelcopter_hyperparameters["max_t"],
    pixelcopter_hyperparameters["gamma"],
    1000,
)
\end{lstlisting}

\paragraph{Публикуем обученную модель на Hub}\label{publish-our-trained-model-on-the-hub-1}

\begin{lstlisting}
repo_id = ""  # TODO Define your repo id {username/Reinforce-{model-id}}
push_to_hub(
    repo_id,
    pixelcopter_policy,  # The model we want to save
    pixelcopter_hyperparameters,  # Hyperparameters
    eval_env,  # Evaluation environment
    video_fps=30
)
\end{lstlisting}

\subsubsection{Дополнительные задания}\label{some-additional-challenges}

Лучший способ научиться --- \textbf{экспериментировать самостоятельно}! Как ты видел, сейчас агент играет не слишком хорошо. Для начала можно увеличить число шагов обучения. Но стоит также подобрать лучшие параметры.

В \href{https://huggingface.co/spaces/huggingface-projects/Deep-Reinforcement-Learning-Leaderboard}{таблице лидеров} можно найти своих агентов. Сможешь подняться на первое место?

Вот несколько идей, которые помогут улучшить позицию в таблице лидеров:

\begin{itemize}
\tightlist
\item
  Обучай агент дольше.
\item
  Пробуй разные гиперпараметры, изучая эксперименты других участников: \url{https://huggingface.co/models?other=reinforce}.
\item
  \textbf{Загружай на Hub новые версии обученных моделей}.
\item
  \textbf{Улучши реализацию для более сложных сред}. Например, что будет, если заменить сеть на свёрточную нейронную сеть (CNN) для обработки изображений кадров в качестве наблюдений?
\end{itemize}

\begin{center}\rule{0.5\linewidth}{0.5pt}\end{center}

\textbf{Поздравляем с завершением модуля}! Материала было немало. И поздравляем с окончанием практики: ты только что с нуля написал своего первого агента глубокого RL на PyTorch и опубликовал его на Hub.

Не бойся продолжать эксперименты с этим модулем, \textbf{улучшая реализацию для более сложных сред}. Например, можно попробовать свёрточную нейронную сеть для обработки кадров-наблюдений.

В следующем модуле мы \textbf{подробнее познакомимся с Unity ML-Agents}, обучая агентов в средах Unity. Это подготовит тебя к \textbf{соревнованиям AI vs AI, где агенты будут противостоять друг другу в снежной битве и футбольном матче}.

Звучит интересно? До встречи в следующем модуле!

Наконец, нам очень хотелось бы \textbf{услышать твоё мнение о курсе и о том, как его можно улучшить}. Если есть предложения, пожалуйста, \href{https://forms.gle/BzKXWzLAGZESGNaE9}{заполни эту форму}.

До встречи в модуле 5!

\paragraph{Продолжай учиться и оставайся молодцом}\label{keep-learning-stay-awesome}

\section{Тест}\label{quiz}

Лучший способ учиться и \href{https://www.coursera.org/lecture/learning-how-to-learn/illusions-of-competence-BuFzf}{не поддаваться иллюзии компетентности} --- \textbf{проверять себя}. Так можно понять, \textbf{какие знания нужно закрепить}.

\subsubsection{Вопрос 1. Каковы преимущества градиента политики перед методами на основе ценности? (Выбери все подходящие варианты)}\label{q1-what-are-the-advantages-of-policy-gradient-over-value-based-methods-check-all-that-apply}

\begin{itemize}
\tightlist
\item
  \textbf{A.} Методы градиента политики могут обучать стохастическую политику.
\item
  \textbf{B.} Методы градиента политики эффективнее при высокой размерности пространства действий и в пространствах непрерывных действий.
\item
  \textbf{C.} Методы градиента политики чаще всего сходятся к глобальному максимуму.
\end{itemize}

\subsubsection{Вопрос 2. Что такое теорема о градиенте политики?}\label{q2-what-is-the-policy-gradient-theorem}

\subsubsection{Вопрос 3. Чем отличаются методы на основе политики от методов градиента политики? (Выбери все подходящие варианты)}\label{q3-whats-the-difference-between-policy-based-methods-and-policy-gradient-methods-check-all-that-apply}

\begin{itemize}
\tightlist
\item
  \textbf{A.} Методы на основе политики являются подмножеством методов градиента политики.
\item
  \textbf{B.} Методы градиента политики являются подмножеством методов на основе политики.
\item
  \textbf{C.} В методах на основе политики параметры \(\theta\) можно оптимизировать \textbf{косвенно}, максимизируя локальную аппроксимацию целевой функции методами восхождения на холм, имитации отжига или эволюционных стратегий.
\item
  \textbf{D.} В методах градиента политики параметры \(\theta\) оптимизируются \textbf{напрямую} с помощью градиентного подъёма целевой функции.
\end{itemize}

\subsubsection{\texorpdfstring{Вопрос 4. Почему для оптимизации \(J(\theta)\) используется градиентный подъём, а не спуск?}{Вопрос 4. Почему используется градиентный подъём, а не спуск?}}\label{q4-why-do-we-use-gradient-ascent-instead-of-gradient-descent-to-optimize-jtheta}

\begin{itemize}
\tightlist
\item
  \textbf{A.} Мы хотим минимизировать \(J(\theta)\), а градиентный подъём задаёт направление наиболее быстрого увеличения \(J(\theta)\).
\item
  \textbf{B.} Мы хотим максимизировать \(J(\theta)\), а градиентный подъём задаёт направление наиболее быстрого увеличения \(J(\theta)\).
\end{itemize}

Поздравляем с завершением теста! Если в каких-то вопросах ошибся, перечитай соответствующие разделы, чтобы закрепить материал.

\subsection{Ответы на тест и пояснения}\label{quiz-solutions-and-explanations}

\subsubsection{Ответ на вопрос 1}\label{q1-solution}

\textbf{Правильные ответы:} A, B.

\textbf{A.} Методы градиента политики могут обучать стохастическую политику.

\textbf{B.} Методы градиента политики эффективнее при высокой размерности пространства действий и в пространствах непрерывных действий.

\textbf{C.} Методы градиента политики чаще всего сходятся к глобальному максимуму. Пояснение: нет, часто градиент политики сходится к локальному максимуму, а не к глобальному оптимуму.

\subsubsection{Ответ на вопрос 2}\label{q2-solution}

\emph{Теорема о градиенте политики} позволяет преобразовать целевую функцию в дифференцируемое выражение, не требующее дифференцирования распределения состояний.

\CourseFigure{policy_gradient_theorem.png}{Теорема о градиенте политики \protect\href{https://huggingface.co/datasets/huggingface-deep-rl-course/course-images/resolve/main/en/unit6/policy_gradient_theorem.png}{[оригинал]}}

\subsubsection{Ответ на вопрос 3}\label{q3-solution}

\textbf{Правильные ответы:} B, C, D.

\textbf{A.} Методы на основе политики являются подмножеством методов градиента политики.

\textbf{B.} Методы градиента политики являются подмножеством методов на основе политики.

\textbf{C.} В методах на основе политики параметры \(\theta\) можно оптимизировать \textbf{косвенно}, максимизируя локальную аппроксимацию целевой функции методами восхождения на холм, имитации отжига или эволюционных стратегий.

\textbf{D.} В методах градиента политики параметры \(\theta\) оптимизируются \textbf{напрямую} с помощью градиентного подъёма целевой функции.

\subsubsection{Ответ на вопрос 4}\label{q4-solution}

\textbf{Правильный ответ:} B.

\textbf{A.} Мы хотим минимизировать \(J(\theta)\), а градиентный подъём задаёт направление наиболее быстрого увеличения \(J(\theta)\).

\textbf{B.} Мы хотим максимизировать \(J(\theta)\), а градиентный подъём задаёт направление наиболее быстрого увеличения \(J(\theta)\).

\section{Заключение}\label{conclusion}

\textbf{Поздравляем с завершением модуля}! Материала было много. И поздравляем с окончанием практики: ты с нуля реализовал своего первого агента глубокого обучения с подкреплением на PyTorch и опубликовал его на Hub.

Продолжай экспериментировать с этим модулем, \textbf{улучшая реализацию для более сложных сред}. Например, почему бы не заменить сеть свёрточной нейронной сетью для обработки кадров-наблюдений?

В следующем модуле мы \textbf{подробнее познакомимся с Unity ML-Agents} и обучим агентов в средах Unity. Благодаря этому ты будешь готов к \textbf{соревнованиям AI vs AI, где агенты противостоят друг другу в снежной битве и футбольном матче}.

Звучит интересно? До встречи в следующий раз!

Наконец, нам очень хотелось бы \textbf{услышать твоё мнение о курсе и о том, как его можно улучшить}. Если есть предложения, пожалуйста, \href{https://forms.gle/BzKXWzLAGZESGNaE9}{заполни эту форму}.

\subsubsection{Продолжай учиться и оставайся молодцом}\label{keep-learning-stay-awesome-1}

\begingroup\small
\section{Глоссарий}\label{glossary}

Этот глоссарий создан сообществом. Мы будем рады твоему вкладу!

\begin{itemize}
\item
  \textbf{Глубокое Q-обучение (Deep Q-Learning):} алгоритм глубокого обучения с подкреплением на основе ценности, использующий глубокую нейронную сеть для аппроксимации Q-значений действий в данном состоянии. Цель --- найти оптимальную политику, максимизирующую ожидаемую суммарную награду, путём обучения ценностей действий.
\item
  \textbf{Методы на основе ценности (value-based):} методы обучения с подкреплением, в которых функция ценности оценивается как промежуточный этап поиска оптимальной политики.
\item
  \textbf{Методы на основе политики (policy-based):} методы обучения с подкреплением, которые непосредственно аппроксимируют оптимальную политику без обучения функции ценности. На практике они выдают распределение вероятностей действий.

  Преимущества методов градиента политики перед методами на основе ценности включают:

  \begin{itemize}
  \tightlist
  \item
    Простоту интеграции: не нужно хранить значения ценности действий.
  \item
    Возможность обучать стохастическую политику: агент исследует пространство состояний, не повторяя одну и ту же траекторию, и избегает проблемы неразличимости состояний по наблюдениям.
  \item
    Эффективность при большой размерности пространства действий и в непрерывных пространствах действий.
  \item
    Более благоприятные свойства сходимости.
  \end{itemize}
\item
  \textbf{Градиент политики (Policy Gradient):} подкласс методов на основе политики, где целью является максимизация качества параметризованной политики с помощью градиентного подъёма. Смысл градиента политики --- управлять распределением вероятностей действий, настраивая политику так, чтобы хорошие действия (повышающие отдачу) выбирались чаще в будущем.
\item
  \textbf{Monte Carlo REINFORCE:} алгоритм градиента политики, который использует оценку отдачи за весь эпизод для обновления параметров политики.
\end{itemize}

Если хочешь помочь улучшить курс, можешь \href{https://github.com/huggingface/deep-rl-class/pulls}{создать Pull Request}.

Создание этого глоссария стало возможным благодаря следующим участникам:

\begin{itemize}
\tightlist
\item
  \href{https://github.com/dcarpintero}{Diego Carpintero}
\end{itemize}

\section{Дополнительная литература}\label{additional-readings}

Эти материалы \textbf{необязательны}: они пригодятся, если захочешь изучить тему глубже.

\subsection{Введение в оптимизацию политики}\label{introduction-to-policy-optimization}

\begin{itemize}
\tightlist
\item
  \href{https://spinningup.openai.com/en/latest/spinningup/rl_intro3.html}{Часть 3. Введение в оптимизацию политики --- документация Spinning Up}.
\end{itemize}

\subsection{Градиент политики}\label{policy-gradient}

\begin{itemize}
\tightlist
\item
  \url{https://johnwlambert.github.io/policy-gradients/}
\item
  \href{https://jonathan-hui.medium.com/rl-policy-gradients-explained-9b13b688b146}{Объяснение градиентов политики в RL}.
\item
  \href{http://incompleteideas.net/book/RLbook2020.pdf}{Глава 13 «Методы градиента политики» в книге Reinforcement Learning: An Introduction Ричарда Саттона и Эндрю Барто}.
\end{itemize}

\subsection{Реализация}\label{implementation}

\begin{itemize}
\tightlist
\item
  \href{https://github.com/pytorch/examples/blob/main/reinforcement_learning/reinforce.py}{Реализация REINFORCE на PyTorch}.
\item
  \href{https://github.com/MrSyee/pg-is-all-you-need}{Реализации алгоритмов от DDPG до PPO}.
\end{itemize}

\endgroup
\end{document}
